% Laget automatisk av elevark/lag_elevark.py fra sim/satellittbaner.html. Endringer her blir overskrevet.
\documentclass[11pt]{article}
\usepackage{elevark}
\begin{document}
\elevarkhode{Fysikk 2 · Gravitasjon}{Satellittbaner}{https://fysikklab.no/sim/satellittbaner.html}{Simuleringen}

En satellitt er hele tiden i fritt fall mot jorda, men den har så stor fart sidelengs at den hele tiden bommer. Med akkurat riktig fart blir banen en sirkel. Med litt mer fart blir den en ellipse, og med nok fart slipper satellitten unna. Hvordan henger farten, høyden og omløpstiden sammen?

\section*{Oppgaver}

\begin{oppgaveliste}
\oppgave \fase{gjett}{Gjett}En satellitt i 400~km høyde får 10~\% større fart enn sirkelbanefarten. Hva skjer med banen? Hva skjer med 10~\% mindre fart?
\svar
\oppgave \fase{test}{Regn}Vis at sirkelbanefarten er $v = \sqrt{GM/r}$. Regn ut farten og omløpstiden for romstasjonen i 400~km høyde.
\svar
\oppgave \fase{test}{Undersøk}Velg geostasjonær bane. Hvor lang er omløpstiden? Hvorfor er det praktisk for TV-satellitter?
\svar
\oppgave \fase{forklar}{Forklar}Se på energigrafen for en sirkelbane og for en ellipse. Hva skjer med $E_{\mathrm{k}}$ og $E_{\mathrm{p}}$ rundt banen, og hva skjer med den samlede energien?
\svar
\oppgave \fase{test}{Regn}Vis at for en sirkelbane er $E_{\mathrm{k}} = -\tfrac{1}{2}E_{\mathrm{p}}$ og $E = \tfrac{1}{2}E_{\mathrm{p}}$. Sjekk med tallene for romstasjonen.
\svar
\oppgave \fase{test}{Undersøk}Hvor stor fart (i prosent av sirkelbanefarten) trengs for at satellitten skal slippe unna? Hvorfor er svaret det samme i alle høyder?
\svar
\oppgave \fase{utfordring}{Utfordring}Romstasjonen bremses litt av den tynne atmosfæren. Hva skjer med høyden og farten? Hvorfor får den større fart når den bremses?
\svar
\end{oppgaveliste}

\end{document}
